% !TEX root = ../metatime_monograph.tex
\chapter{No-Go Theorem for a Common Up-Sector Baseline}
\label{app:up_sector_common_baseline_no_go}


% CITATION_CONTEXT_V10_9_BEGIN
\paragraph{Established context.}
The no-go result is an internal algebraic theorem for the frozen v10.2 trace.
Its scientific handling follows falsification and preproducibility principles
\cite{Popper1959,Lakatos1978,Nosek2018,Stark2018}: an excluded architecture is
retained as excluded rather than rescued by an unregistered correction.
% CITATION_CONTEXT_V10_9_END

\section{Scope and status}

Appendix~\ref{app:sector_holonomy_release} isolated a large first-generation
up-quark residual while the charm and top channels remained much closer to their
validation scales.  The v10.3 preregistration proposed testing a common additive
baseline for the full up-type sector.  This appendix proves that such a baseline
cannot close the fixed v10.2 trace.

The status is

\[
\boxed{
\text{technical PASS}
\;\land\;
\text{common-baseline architecture FAIL}
\;\land\;
\text{physical spectrum OPEN}
}
\]

The calculation is retrospective because the validation residuals have already
been inspected.  It proposes no new mass coefficient and permits no canonical
promotion.

\section{Common-baseline architecture}

Suppose the up-type masses are represented by

\begin{equation}
\ln\frac{m_{u_g}}{m_e}
=
B_{\rm up}+\Delta_{\rm up}(g),
\qquad g\in\{u,c,t\},
\label{eq:common_up_baseline}
\end{equation}

where the same dimensionless baseline $B_{\rm up}$ acts on all three channels
and the relative trace $\Delta_{\rm up}(g)$ is held fixed.

Let

\begin{equation}
r_g=\ln\frac{m_g^{\rm pred}}{m_g^{\rm val}}
\label{eq:up_residual_definition}
\end{equation}

be the logarithmic residual.  Adding a common baseline shift $B$ gives

\begin{equation}
r'_g=r_g+B.
\label{eq:translated_up_residual}
\end{equation}

\section{Translation-invariance theorem}

\begin{theorem}[Common-baseline no-go]
For a fixed residual vector $\{r_u,r_c,r_t\}$, a common additive shift cannot
alter any pairwise residual difference.  A common shift exists such that
$|r'_g|\leq\varepsilon$ for every up-type channel if and only if

\[
\max_g r_g-\min_g r_g\leq 2\varepsilon.
\]

The smallest achievable uniform absolute logarithmic error is

\[
\varepsilon_{\min}
=
\frac{\max_g r_g-\min_g r_g}{2}.
\]
\end{theorem}

\begin{proof}
For any pair $g,h$,

\[
r'_g-r'_h=(r_g+B)-(r_h+B)=r_g-r_h.
\]

Thus the residual spread is invariant.  The condition $|r_g+B|\leq\varepsilon$
requires

\[
B\in[-r_g-\varepsilon,-r_g+\varepsilon]
\]

for every $g$.  These intervals have a common intersection exactly when the
largest lower endpoint does not exceed the smallest upper endpoint, which is
equivalent to $\max r-\min r\leq2\varepsilon$.  The minimax translation centres
the extreme residuals around zero, giving half the spread.  \qed
\end{proof}

\section{Application to the v10.2 trace}

The already inspected v10.2 residuals are

\[
r_u=3.1857238256,
\qquad
r_c=0.0818052322,
\qquad
r_t=0.2293790515.
\]

Their invariant spread is

\begin{equation}
\Delta r
=
\max r-\min r
=
3.1039185934.
\label{eq:up_residual_spread}
\end{equation}

Therefore

\begin{equation}
\boxed{
\varepsilon_{\min}=1.5519592967
}
\label{eq:up_minimax_error}
\end{equation}

and at least one channel must retain a multiplicative error not smaller than

\begin{equation}
\boxed{
\exp(\varepsilon_{\min})=4.7207104.
}
\label{eq:up_minimax_factor}
\end{equation}

The optimal minimax shift is

\[
B_*=-\frac{\max r+\min r}{2}=-1.6337645289,
\]

for which

\[
|r'_u|=|r'_c|=1.5519592967,
\qquad
|r'_t|=1.4043854774.
\]

Hence no common baseline can bring all three up-type channels into the
approximately $0.254$ logarithmic envelope achieved by the other v10.2
non-anchor channels.

\section{Interpretation}

The no-go theorem applies to the architecture
$B_{\rm up}+\Delta_{\rm up}(g)$ with the v10.2 relative trace frozen.  It does
not disprove the Collatz quarter-power exponent and does not exclude a richer
up-sector geometry.

The repository already distinguishes the structural labels
\path{light_quark_seed} and \path{heavy_quark_resonance}.  A future
candidate may investigate whether these pre-existing sectors require distinct
structural baselines, or whether the relative up-sector release itself must be
rederived.  Such work must preserve the quarantine of the archived heavy-sector
bridge and may not introduce a correction keyed directly to the particle name
$u$.

Any numerical comparison using the already inspected charged-fermion mass table
remains retrospective.  Independent preregistered validation is still required
for promotion.

The executable audit is located at
\path{TIR/validation/up_sector_common_baseline_no_go_v10_4.py}.
